\documentclass[a4paper]{exam}
\usepackage[svgnames]{xcolor}
\usepackage{amsmath,amssymb,amsthm,titlesec,minted}

\newcommand
\lralign[2]{%
  \noindent\makebox[\textwidth]{\(\displaystyle#1\)\hfill\ifx&#2&\else(#2)\fi}\vspace{2ex}
}

\renewcommand{\thesubsection}{\alph{subsection}.}
\renewcommand{\thesubsubsection}{\roman{subsubsection}.}

\definecolor{faintgreen}{rgb}{0.5,1,0.5}
\definecolor{faintred}{rgb}{1,0.5,0.5}
\definecolor{faintblue}{rgb}{0.5,0.5,1}
\definecolor{faintyellow}{rgb}{1,1,0.5}
\definecolor{codebg}{RGB}{40,42,54}

\titleformat{\section}{\normalfont\Large\bfseries}{}{0pt}{}
\titleformat{\subsection}{\normalfont\large\bfseries}{}{0pt}{}
\titleformat{\subsubsection}{\normalfont\normalsize\bfseries\itshape}{}{0pt}{}
\titlespacing*{\section}{1em}{1.6\baselineskip}{0.6\baselineskip}
\titlespacing*{\subsection}{2em}{1.2\baselineskip}{0.4\baselineskip}
\titlespacing*{\subsubsection}{4em}{1.0\baselineskip}{0.3\baselineskip}

\setminted{style=dracula,bgcolor=codebg,breaklines,autogobble,fontsize=\small,tabsize=4,baselinestretch=1}
\NewDocumentCommand{\codefromfile}{O{}mm}{\par\vspace{0.4em}\inputminted[#1]{#2}{#3}\par\vspace{0.4em}}
\NewDocumentEnvironment{Section}{m}
{%
  \section*{#1}%
  \begingroup
  \leftskip=1em
}
{%
  \par
  \endgroup
}

\NewDocumentEnvironment{Subsection}{m}
{%
  \subsection*{#1}%
  \begingroup
  \leftskip=2em
}
{%
  \par
  \endgroup
}

\setlength{\parindent}{0in}
\setlength{\oddsidemargin}{0in}
\setlength{\textwidth}{6.5in}
\setlength{\textheight}{8.8in}
\setlength{\topmargin}{0in}
\setlength{\headheight}{18pt}

\printanswers
\title{Week 6}
\author{Radhesh Goel}
\date{\today}

\begin{document}
\maketitle

\section*{Question 1.}
Rewrite Equations (2)--(4) in a form suitable for Jacobi iteration.
\begin{solution}
  The discretised heat equations are
  \[
    \begin{aligned}
      -k\left( \frac{T_L - 2T_0 + T_1}{\Delta x^2} \right) & = q \\[0.5em]
      -k\left( \frac{T_0 - 2T_1 + T_2}{\Delta x^2} \right) & = q \\[0.5em]
      -k\left( \frac{T_1 - 2T_2 + T_R}{\Delta x^2} \right) & = q
    \end{aligned}
  \]

  Define
  \[
    C = \frac{q\Delta x^2}{k}
  \]

  Using
  \[
    k = 170, \qquad q = 4 \times 10^5, \qquad \Delta x = 0.1
  \]

  gives
  \[
    C = \frac{(4\times10^5)(0.1)^2}{170}
      = \frac{400}{17}
      \approx 23.5294\ \mathrm{K}
  \]

  Multiplying each equation by
  \[
    -\frac{\Delta x^2}{k}
  \]

  gives
  \[
    \begin{aligned}
      T_L - 2T_0 + T_1 & = -C \\
      T_0 - 2T_1 + T_2 & = -C \\
      T_1 - 2T_2 + T_R & = -C
    \end{aligned}
  \]

  Rearranging for each unknown temperature,
  \[
    \begin{aligned}
      T_0 & = \frac{T_L + T_1 + C}{2} \\
      T_1 & = \frac{T_0 + T_2 + C}{2} \\
      T_2 & = \frac{T_1 + T_R + C}{2}
    \end{aligned}
  \]

  Therefore, the Jacobi iteration equations are
  \[
    \boxed{\begin{aligned}
            T_0^{(n+1)} &= \frac{T_L + T_1^{(n)} + C}{2} \\[0.4em]
            T_1^{(n+1)} &= \frac{T_0^{(n)} + T_2^{(n)} + C}{2} \\[0.4em]
            T_2^{(n+1)} &= \frac{T_1^{(n)} + T_R + C}{2}
        \end{aligned}}
  \]

  Substituting the boundary temperatures and the value of \(C\),
  \[
    \boxed{\begin{aligned}
        T_0^{(n+1)} &= 161.7647 + 0.5T_1^{(n)}, \\[0.4em]
        T_1^{(n+1)} &= 11.7647 + 0.5T_0^{(n)} + 0.5T_2^{(n)}, \\[0.4em]
        T_2^{(n+1)} &= 386.7647 + 0.5T_1^{(n)}.
        \end{aligned}}
  \]
\end{solution}
\end{document}
